\documentclass[aspectratio=169, xcolor=dvipsnames]{beamer}

\useoutertheme{miniframes}
\useinnertheme{circles}
\usepackage{appendixnumberbeamer}
\setbeamertemplate{footline}[frame number]

\definecolor{main_color}{HTML}{7D9D33}
\definecolor{secondary_color}{HTML}{CED38C}

\setbeamercolor{palette primary}{bg=secondary_color,fg=white}
\setbeamercolor{palette secondary}{bg=secondary_color,fg=white}
\setbeamercolor{palette tertiary}{bg=secondary_color,fg=white}
\setbeamercolor{palette quaternary}{bg=secondary_color,fg=white}
\setbeamercolor{structure}{fg=main_color}
\setbeamercolor{section in toc}{fg=secondary_color}
\setbeamercolor{subsection in head/foot}{bg=main_color,fg=white}

\usepackage{booktabs}
\usepackage{graphicx}
\usepackage{subcaption}
\usepackage{caption}

% Set the path for figures
\graphicspath{{/Users/ohouck/OneDrive - The University of Chicago/ai_weather_ag/figures/}}

\title{Tailoring Machine Learning Weather Predictions for Local Impacts}
\author{\textbf{Ozma Houck, James Franke}}
\date{}

\usepackage{bbm}

\begin{document}

\maketitle


\begin{frame}{Research Goal and Overview}
\begin{itemize}
    \item Show the feasibility of finetuning (localizing) weather forecasts for specific outcomes using minimal compute power
    \item Demonstrate the method using 2 meter temperature and 10 meter wind speed across a variety of regions and lead times
    \item Show that an ML forecast like Pangu can be corrected in the same way as traditional NWP forecast like IFS
\end{itemize}

\vspace{0.5em}
\textbf{Key Contribution:} Demonstrating a simple method to significantly improve forecast accuracy using limited computational power.
\end{frame}


\begin{frame}{Method}
\begin{itemize}
    \item Uses Pangu and IFS HRES forecasts with lead times from 1 to 7 days
    \begin{itemize}
        \item Target Pangu predictions using ERA5 and IFS HRES using its t=0 initialization
    \end{itemize}
    \item Filter to a single variable and lead time, in a particular region
    \item Train simple MLP network that takes in Pangu/IFS forecast outputs and minimizes the forecast error
    \begin{itemize}
        \item Train using a single forecast per day for years 2018-2021
        \item Test using daily forecasts from 2022
    \end{itemize}
\end{itemize}
\end{frame}


\begin{frame}{Headline Result: Simple MLP Seems Effective at Correcting Forecasts}
\begin{itemize}
    \item Simple MLP method works across variety of regions and climates
    \item Training a correction model for a 10x10 degree region takes about 2.5 minutes on an M3 max MacBook Pro
    \item Similar improvements to both Pangu and IFS forecasts
\end{itemize}

\begin{figure}[ht]
    \centering
    \begin{subfigure}[t]{0.48\textwidth}
        \centering
        \includegraphics[width=\linewidth]{finetuning/pangu/comparison/mse_comparison_pangu_trained_with_2m_temperature_output2m_temperature_mlp512x5.png}
        \caption{2m Temperature}
    \end{subfigure}
    \hfill
    \begin{subfigure}[t]{0.48\textwidth}
        \centering
        \includegraphics[width=\linewidth]{finetuning/pangu/comparison/mse_comparison_pangu_trained_with_10m_wind_speed_output10m_wind_speed_mlp512x5.png}
        \caption{Wind Speed}
    \end{subfigure}
\end{figure}
\end{frame}

\begin{frame}{See Similar Percent Improvements Across Different Lead Times}
\begin{figure}[ht]
    \centering
    \begin{subfigure}[t]{0.48\textwidth}
        \centering
        \includegraphics[width=\linewidth]{finetuning/pangu/lead_time/india/10x10/leadtime_rmse_pct_2m_temperature_trainedwith_2m_temperature_mlp512x5.png}
        \caption{2m Temperature: 10x10 degree}
    \end{subfigure}
    \hfill
    \begin{subfigure}[t]{0.48\textwidth}
        \centering
        \includegraphics[width=\linewidth]{finetuning/pangu/lead_time/india/10x10/leadtime_rmse_pct_10m_wind_speed_trainedwith_10m_wind_speed_mlp512x5.png}
        \caption{Wind Speed: 10x10 degree}
    \end{subfigure}
\end{figure}
\end{frame}

\begin{frame}{How to Explore Heterogeneity by Climate Region}
\begin{itemize}
    \item Define tropical, temperate, and arid pixels on an 0.25 degree grid using broad Koppen-Geiger classes
    \item Train fine-tuning model on randomly sampled 2x2 degree patches
\end{itemize}

\begin{figure}
    \centering
    \includegraphics[width=0.7\linewidth]{finetuning/climate_zones_with_patches.png}
\end{figure}
\end{frame}

\begin{frame}{Heterogeneity by Climate Region: 2m Temperature}
    \begin{figure}[ht]
    \centering
    % Top row with 2 figures
    \begin{subfigure}[t]{0.45\textwidth}
    \centering
    \includegraphics[width=\linewidth]{finetuning/pangu/lead_time/temperate/2x2/leadtime_rmse_pct_2m_temperature_trainedwith_2m_temperature_mlp512x5_bootstrap.png}
    \end{subfigure}
    \hfill
    \begin{subfigure}[t]{0.45\textwidth}
    \centering
    \includegraphics[width=\linewidth]{finetuning/pangu/lead_time/arid/2x2/leadtime_rmse_pct_2m_temperature_trainedwith_2m_temperature_mlp512x5_bootstrap.png}
    \end{subfigure}
    
    \vspace{0.5em} % Add vertical space between rows
    
    % Bottom row with 1 centered figure
    \begin{subfigure}[t]{0.45\textwidth}
    \centering
    \includegraphics[width=\linewidth]{finetuning/pangu/lead_time/tropical/2x2/leadtime_rmse_pct_2m_temperature_trainedwith_2m_temperature_mlp512x5_bootstrap.png}
    \end{subfigure}
    \end{figure}
    \end{frame}
    
    \begin{frame}{Heterogeneity by Climate Region: 10m Wind Speed}
    \begin{figure}[ht]
    \centering
    % Top row with 2 figures
    \begin{subfigure}[t]{0.45\textwidth}
    \centering
    \includegraphics[width=\linewidth]{finetuning/pangu/lead_time/temperate/2x2/leadtime_rmse_pct_10m_wind_speed_trainedwith_10m_wind_speed_mlp512x5_bootstrap.png}
    \end{subfigure}
    \hfill
    \begin{subfigure}[t]{0.45\textwidth}
    \centering
    \includegraphics[width=\linewidth]{finetuning/pangu/lead_time/arid/2x2/leadtime_rmse_pct_10m_wind_speed_trainedwith_10m_wind_speed_mlp512x5_bootstrap.png}
    \end{subfigure}
    
    \vspace{0.5em} % Add vertical space between rows
    
    % Bottom row with 1 centered figure
    \begin{subfigure}[t]{0.45\textwidth}
    \centering
    \includegraphics[width=\linewidth]{finetuning/pangu/lead_time/tropical/2x2/leadtime_rmse_pct_10m_wind_speed_trainedwith_10m_wind_speed_mlp512x5_bootstrap.png}
    \end{subfigure}
    \end{figure}
    \end{frame}



\begin{frame}{Amazon}
\begin{figure}[ht]
    \centering
    \begin{subfigure}[t]{0.48\textwidth}
        \centering
        \includegraphics[width=\linewidth]{finetuning/pangu/lead_time/amazon/10x10/leadtime_rmse_pct_2m_temperature_trainedwith_2m_temperature_mlp512x5.png}
        \caption{2m Temperature: 10x10 degree}
    \end{subfigure}
    \hfill
    \begin{subfigure}[t]{0.48\textwidth}
        \centering
        \includegraphics[width=\linewidth]{finetuning/pangu/lead_time/amazon/10x10/leadtime_rmse_pct_10m_wind_speed_trainedwith_10m_wind_speed_mlp512x5.png}
        \caption{Wind Speed: 10x10 degree}
    \end{subfigure}
\end{figure}
\end{frame}

\begin{frame}{US South}
\begin{figure}[ht]
    \centering
    \begin{subfigure}[t]{0.48\textwidth}
        \centering
        \includegraphics[width=\linewidth]{finetuning/pangu/lead_time/usa_south/10x10/leadtime_rmse_pct_2m_temperature_trainedwith_2m_temperature_mlp512x5.png}
        \caption{2m Temperature: 10x10 degree}
    \end{subfigure}
    \hfill
    \begin{subfigure}[t]{0.48\textwidth}
        \centering
        \includegraphics[width=\linewidth]{finetuning/pangu/lead_time/usa_south/10x10/leadtime_rmse_pct_10m_wind_speed_trainedwith_10m_wind_speed_mlp512x5.png}
        \caption{Wind Speed: 10x10 degree}
    \end{subfigure}
\end{figure}
\end{frame}

\begin{frame}{British Columbia}
\begin{figure}[ht]
    \centering
    \begin{subfigure}[t]{0.48\textwidth}
        \centering
        \includegraphics[width=\linewidth]{finetuning/pangu/lead_time/british_columbia/10x10/leadtime_rmse_pct_2m_temperature_trainedwith_2m_temperature_mlp512x5.png}
        \caption{2m Temperature: 10x10 degree}
    \end{subfigure}
    \hfill
    \begin{subfigure}[t]{0.48\textwidth}
        \centering
        \includegraphics[width=\linewidth]{finetuning/pangu/lead_time/british_columbia/10x10/leadtime_rmse_pct_10m_wind_speed_trainedwith_10m_wind_speed_mlp512x5.png}
        \caption{Wind Speed: 10x10 degree}
    \end{subfigure}
\end{figure}
\end{frame}

\begin{frame}{Previous Paper Corrected IFS Used a U-net}
\begin{itemize}
    \item \cite{hanDeepLearningMethod2021} trained a U-Net to correct IFS Forecasts in Northern China using 12 years of data and a NVIDIA GPU
    \begin{itemize}
        \item Outperforms traditional method of correcting by subtracting mean forecast error
    \end{itemize}
    \item Uses RMSE error instead of MSE but I find a larger percent improvement for both 2m temperature and 10m wind speed
\end{itemize}
\end{frame}

\begin{frame}{What I'm Actively Working on}
\begin{itemize}
    \item Rewriting the databuild to be able to take advantage of the new WeatherbenchX package and also be able to run on a server and be more flexible
    \item Implementing a U-Net model for fine-tuning
    \begin{itemize}
        \item Early results suggest it significantly ups compute time without much improved performance
    \end{itemize}
    \item Implement ANO Correction (subtracting mean error) method to compare to MLP correction
\end{itemize}
\end{frame}

\end{document}